% ===========================================================================
%  sviluppatore/04-motore/02-decisione.tex
% ===========================================================================
\section{Decisione ed esecuzione sono due istanti}\label{sez:decisione-esecuzione}

È la regola più importante del motore, e si riduce a due righe di codice.

\begin{codiceInterno}{python}{Le due righe che tengono in piedi tutto}{backtest/engine.py:740}
            # t_execution = t (open T)
            # t_decision = data point precedente in all_dates (close T-1)
            t_execution = t
            t_decision = all_dates[i - 1]
\end{codiceInterno}

\figuraTikz{decisione-esecuzione}
  {I due istanti. Tutto ciò che sta nella fascia a destra non esiste ancora
   quando si decide.}
  {decisione-esecuzione}

\subsection{Perché due istanti e non uno}

Una strategia che decide alla chiusura di $T$ e compra alla chiusura di $T$
non è implementabile: nel momento in cui si conosce il prezzo di chiusura, la
sessione è finita. Le implementazioni ingenue lo fanno lo stesso, e il
backtest che ne esce è eccellente.

Qui la decisione avviene con i dati disponibili fino alla chiusura di $T-1$ e
l'ordine viene eseguito al prezzo di apertura di $T$:

\begin{codiceInterno}{python}{La fase di esecuzione}{backtest/engine.py:945}
            # ---- EXECUTION PHASE: open T+1 = t_execution ----
            # Per ogni target ottengo l'open price a t_execution.
            ...
                df_sym = data.get(sym)
                if df_sym is None or t_execution not in df_sym.index:
                    # Symbol non disponibile, skip silenzioso
                    continue
                open_price = float(df_sym.loc[t_execution, "open"])
\end{codiceInterno}

\subsection{La firma che non si può cambiare}

Le strategie ricevono i due istanti come argomenti \emph{nominati}:

\begin{codiceInterno}{python}{Il contratto}{backtest/engine.py:834}
                        t_decision=t_decision,
                        t_execution=t_execution,
\end{codiceInterno}

\noindent
Il motore li passa per nome, quindi una strategia che chiami i suoi parametri
in modo diverso fa fermare il backtest alla prima barra con un
\cod{TypeError}. Il messaggio di errore del motore lo dice esplicitamente,
ricordando la firma richiesta:

\begin{codiceInterno}{python}{}{backtest/engine.py:857}
                            "Firma richiesta: generate_signals(self, t_decision, "
                            "t_execution, indicators_by_symbol, portfolio) -> "
\end{codiceInterno}

\begin{attenzione}{Il bug che migliora i risultati}
Dentro \simbolo{generate\_signals} si legge solo fino a \cod{t\_decision}:

\begin{center}
\begin{tabular}{@{}L{62mm}L{76mm}@{}}
\toprule
\textbf{Corretto} & \textbf{Lookahead} \\
\midrule
\cod{ind.close.loc[:t\_decision]} & \cod{ind.close.loc[:t\_execution]} \\
\cod{serie.shift(1)} prima di decidere & \cod{serie} senza shift \\
\bottomrule
\end{tabular}
\end{center}

\noindent
La seconda colonna non dà errore, non produce avvisi, e fa salire lo Sharpe.
È per questo che la separazione è imposta dalla firma invece che raccomandata
in un commento: una regola che si può violare senza accorgersene viene
violata.
\end{attenzione}

\subsection{La stessa regola altrove}

La causalità non riguarda solo le strategie. Si ripresenta identica in tre
altri punti, e in ciascuno è scritta nel codice come vincolo:

\begin{center}
\begin{tabular}{@{}L{34mm}L{104mm}@{}}
\toprule
\textbf{Dove} & \textbf{Come} \\
\midrule
\rigaTenue
Regime & Lo snapshot è calcolato con \cod{as\_of=True}, che usa solo dati
         anteriori a \cod{t\_decision}. \\
Overlay & Il moltiplicatore si calcola dal PnL \emph{base} e si applica
          shiftato: «la decisione di oggi non può usare il rendimento di
          oggi». \\
\rigaTenue
Regole di regime & Niente statistiche sull'intero campione: un quantile
          calcolato su tutta la storia è lookahead travestito da soglia.
          Gli helper condivisi (\cod{\_expanding\_median},
          \cod{\_hysteresis}) esistono per non doverlo riscrivere ogni volta. \\
Normalizzazione della volatilità & La vol trailing viene shiftata di un
          giorno nel backtest (\vediSez{ensemble}). \\
\bottomrule
\end{tabular}
\end{center}
